\section{Introduction}

The rapid evolution of foundation models is expanding the boundary of what can be natively learned and performed within a single model.
Many capabilities are progressively internalized from external systems into the model itself, such as reasoning~\cite{guo2025deepseek,jaech2024openai} and multimodality~\cite{xu2026deepseek,team2024gemini}.
Memory represents another fundamental capability, enabling models to accumulate experience and adapt over extended interactions~\cite{zhang2025survey,hu2025memory}.
Our previous work, Metis~\cite{zhang2026metis}, first introduced memory foundation models to internalize memory as a native capability of foundation models.
This offers a path beyond external context and enables models to continuously accumulate information and experience, which marks a critical step towards memory-native evolution.

To improve memory capabilities, we introduce \textbf{MemT-1.0}, a new architecture with substantial upgrades over Metis, while retaining the structure of dynamic local memory and trainable hyper memory.
It replaces token selection with a stable token-level update mechanism based on Gated DeltaNet-2 (GDN-2)~\cite{hatamizadeh2026gated}.
To further scale native memory capability, we introduce Mixture-of-Storage (MoS) to expand the capacity of native memory.
We release a model family comprising \textbf{MemT-1.0-Nano} with 36B parameters (3B activated, 10M dynamic memory) and \textbf{MemT-1.0-Flash} with 288B parameters (17B activated, 336M dynamic memory).

Following a deliberate backbone preparation, we develop a two-stage training pipeline for MemT-1.0.
The first stage performs continued pre-training (CPT) to establish language modeling capabilities with native memory~\cite{gururangan2020don}.
The post-training stage further adapts the model to downstream memory tasks while preserving its general capabilities~\cite{ouyang2022training}.
To support this training pipeline, we construct a large-scale dataset of 34B tokens.
Across 11 memory benchmarks, our models outperform parametric memory baselines on average, showing effective utilization of historical information without accessing the original context during inference.
They also largely preserve general capabilities under the noise setting.
Our evaluation shows an increasing end-to-end inference latency and storage overhead advantage over full context on longer histories.

To support training and deployment, we develop a distributed infrastructure, which preserves memory isolation across independent trajectories.
We organize stateful inference around separate memory storage and utilization procedures, with shared model parameters and trajectory-specific memory states.
For a fixed memory configuration, the persistent state size remains constant as interactions accumulate.
In addition, we explore the scaling behavior of MemT-1.0.
Unlike general foundation models~\cite{kaplan2020scalinglawsneurallanguage}, memory foundation models introduce memory capacity as a critical scaling dimension, so we conduct scaling studies over memory capacity and training data.
These findings provide valuable guidance for building stronger memory foundation models.

We further apply MemT-1.0 to experience-guided coding agents, where retrieved procedural experience is internalized into native memory to guide subsequent actions.
On SWE-bench Verified~\cite{chowdhury2024swebenchverified} and SWE-bench Multilingual~\cite{Yang2025SWEsmithSD}, our method improves average task success and consistency across repeated runs over agents without experience, while achieving performance comparable to explicit textual experience.
These findings extend the role of native memory from information recall to procedural guidance in multi-step execution.

Our main contributions are summarized as follows: \\
$\bullet$ We introduce MemT-1.0, a memory foundation model architecture with improved memory storage and utilization mechanisms, together with MoS for scalable memory capacity.\\
$\bullet$ Following a deliberate backbone preparation, we design a two-stage training pipeline including CPT and post-training stages, with a corresponding infrastructure design.
We construct a training corpus of 34B tokens, and systematically evaluate MemT-1.0 on both memory and general benchmarks. The results show that MemT-1.0 achieves superior memory performance while largely preserving general capabilities. \\
$\bullet$ We explore the scaling behavior of memory capabilities in MemT-1.0.
We further deploy MemT-1.0 in practical agentic coding tasks to demonstrate its effectiveness and potential in enabling memory-native evolution.
We also discuss its current limitations and promising directions for future development.
